% Generated draft: requires mathematical and source review.
\documentclass[11pt]{article}
\usepackage[margin=25mm]{geometry}
\usepackage{fontspec}
\setmainfont{Latin Modern Roman}
\usepackage{amsmath,amssymb,mathtools,unicode-math}
\setmathfont{Latin Modern Math}
\usepackage{enumitem,needspace,xurl,hyperref}
\setlength{\parindent}{0pt}
\setlength{\parskip}{5pt}
\setlength{\emergencystretch}{3em}
\newcommand{\N}{\mathbb{N}}
\newcommand{\Z}{\mathbb{Z}}
\newcommand{\Q}{\mathbb{Q}}
\newcommand{\R}{\mathbb{R}}
\newcommand{\C}{\mathbb{C}}
\newcommand{\F}{\mathbb{F}}
\newcommand{\A}{\mathbb{A}}
\newcommand{\PP}{\mathbb{P}}
\newcommand{\E}{\mathbb{E}}
\newcommand{\eps}{\varepsilon}
\newcommand{\dd}{\,\mathrm{d}}
\begin{document}
\section*{UnsolvedMath ID: 30000932}
\textbf{Source problem code:} catalogue entry
\subsection{Definitions and mathematical statement}
The source considers \[u_t^\varepsilon+\partial_x f(u^\varepsilon)=\partial_x R(\varepsilon u_x^\varepsilon,\varepsilon^2u_{xx}^\varepsilon,\ldots),\qquad \varepsilon>0,\] and its proposed sharp-interface limit $u_t+\partial_x f(u)=0$. The state space, flux $f$, regularization $R$, and the derivatives hidden in the ellipsis are not specified. In a one-dimensional Riemann formulation, the initial data would be $u(0,x)=u_L$ for $x<0$ and $u(0,x)=u_R$ for $x>0$, but the source supplies no states $u_L,u_R$ or admissible state domain. A discontinuity from $u_-$ to $u_+$ moving with speed $\sigma$ must satisfy the Rankine--Hugoniot relation \[\sigma(u_+-u_-)=f(u_+)-f(u_-).\] For a specified convex entropy pair $(\eta,q)$, entropy admissibility requires $\partial_t\eta(u)+\partial_xq(u)\leq0$ in distributions, giving $q(u_+)-q(u_-)-\sigma(\eta(u_+)-\eta(u_-))\leq0$ at the jump. The original explicitly asks for a kinetic relation in addition to this entropy inequality; that qualification is retained although abbreviated in the clean/current version.\par In a strictly hyperbolic setting, an undercompressive shock has fewer incoming characteristics than a Lax compressive shock; its precise characteristic count requires the system and shock family. A kinetic relation is additional model-dependent information constraining the jump states and speed, expressible schematically as $\mathcal K(u_-,u_+,\sigma)=0$ or a specified relation between these data. It is not presumed to be a single-valued function. The source asks which jumps are selected by the singular perturbation, how undercompressive shocks are characterized, and what kinetic relation produces the corresponding unique admissible Riemann solution. Traveling profiles $u^\varepsilon(t,x)=U((x-\sigma t)/\varepsilon)$ with $U(\pm\infty)=u_\pm$ are a usual way of examining such selection; no equivalence between existence of a profile and all vanishing-perturbation limits is asserted without the missing model hypotheses. In particular $f$ and an entropy inequality alone do not specify the source's regularization-dependent kinetic relation. The 2026-08-21 triage labels this open as extracted, says model and data are missing, and reports classified solution families for related models without verifying this exact claim as solved.
\subsection{Short English statement}
Determine the shocks selected by the stated singular regularization and the kinetic condition supplementing entropy admissibility for a unique Riemann solver. The flux, regularization, states, and scalar-versus-system setting are not supplied, so no particular kinetic relation can be recovered.
\subsection{Sources}
\begin{itemize}\item Local UnsolvedMath registry, permanent ID 30000932, source problem catalogue entry, \emph{Admissibility and Kinetic Relations for Undercompressive Shocks}, proposed year 2008; no individual proposer supplied.\item \emph{Mini-Workshop: Hyperbolic Aspects of Phase Transition Dynamics} (2008), Oberwolfach Reports, DOI 10.4171/owr/2008/11, pp.\ 513--556: \url{https://doi.org/10.4171/owr/2008/11}. Supplied literature label: original OWR report.\item Supplied literature label: nonclassical Riemann solver, \url{https://doi.org/10.1017/S030821050000158X}.\end{itemize}
\end{document}
